% Figure for Classical Mechanics §B10.3 (effective potential of the heavy symmetric top, units I = Mgl = 1, L3 = I3 omega3 = 4, Lz = L3 cos 60 deg = 2 (released from rest at theta0 = 60 deg), E' = Mgl cos theta0 = 0.5; turning points 60.0 and 66.8 deg, minimum at 63.3 deg).
\begin{tikzpicture}[font=\small]
  \begin{axis}[nfig, width=9cm, height=6cm, xmin=50, xmax=78, ymin=0.42, ymax=0.95,
      xlabel={$\theta$ (degrees)}, ylabel={$U_{\text{eff}}(\theta)/Mg\ell$}, xtick={50,55,60,65,70,75}, ytick={0.5,0.6,0.7,0.8,0.9}, clip=false]
    \addplot[figB] coordinates {(50.30,0.8990) (50.44,0.8893) (50.57,0.8797) (50.71,0.8703) (50.85,0.8610) (50.99,0.8518) (51.12,0.8427) (51.26,0.8338) (51.40,0.8250) (51.53,0.8164) (51.67,0.8079) (51.81,0.7995) (51.95,0.7912) (52.08,0.7831) (52.22,0.7751) (52.36,0.7672) (52.49,0.7594) (52.63,0.7517) (52.77,0.7442) (52.91,0.7368) (53.04,0.7295) (53.18,0.7224) (53.32,0.7153) (53.46,0.7084) (53.59,0.7016) (53.73,0.6949) (53.87,0.6883) (54.00,0.6818) (54.14,0.6754) (54.28,0.6692) (54.42,0.6630) (54.55,0.6570) (54.69,0.6511) (54.83,0.6453) (54.96,0.6396) (55.10,0.6340) (55.24,0.6285) (55.38,0.6231) (55.51,0.6179) (55.65,0.6127) (55.79,0.6076) (55.92,0.6027) (56.06,0.5978) (56.20,0.5930) (56.34,0.5884) (56.47,0.5838) (56.61,0.5794) (56.75,0.5750) (56.88,0.5708) (57.02,0.5666) (57.16,0.5626) (57.30,0.5586) (57.43,0.5548) (57.57,0.5510) (57.71,0.5474) (57.85,0.5438) (57.98,0.5403) (58.12,0.5369) (58.26,0.5337) (58.39,0.5305) (58.53,0.5274) (58.67,0.5244) (58.81,0.5215) (58.94,0.5186) (59.08,0.5159) (59.22,0.5133) (59.35,0.5108) (59.49,0.5083) (59.63,0.5059) (59.77,0.5037) (59.90,0.5015) (60.04,0.4994) (60.18,0.4974) (60.31,0.4955) (60.45,0.4937) (60.59,0.4919) (60.73,0.4903) (60.86,0.4887) (61.00,0.4872) (61.14,0.4858) (61.27,0.4845) (61.41,0.4833) (61.55,0.4822) (61.69,0.4811) (61.82,0.4802) (61.96,0.4793) (62.10,0.4785) (62.24,0.4778) (62.37,0.4771) (62.51,0.4766) (62.65,0.4761) (62.78,0.4758) (62.92,0.4755) (63.06,0.4752) (63.20,0.4751) (63.33,0.4751) (63.47,0.4751) (63.61,0.4752) (63.74,0.4754) (63.88,0.4757) (64.02,0.4760) (64.16,0.4765) (64.29,0.4770) (64.43,0.4776) (64.57,0.4783) (64.70,0.4790) (64.84,0.4799) (64.98,0.4808) (65.12,0.4818) (65.25,0.4829) (65.39,0.4840) (65.53,0.4853) (65.66,0.4866) (65.80,0.4880) (65.94,0.4894) (66.08,0.4910) (66.21,0.4926) (66.35,0.4943) (66.49,0.4961) (66.63,0.4980) (66.76,0.4999) (66.90,0.5019) (67.04,0.5040) (67.17,0.5062) (67.31,0.5085) (67.45,0.5108) (67.59,0.5132) (67.72,0.5157) (67.86,0.5182) (68.00,0.5209) (68.13,0.5236) (68.27,0.5264) (68.41,0.5292) (68.55,0.5322) (68.68,0.5352) (68.82,0.5383) (68.96,0.5415) (69.09,0.5447) (69.23,0.5481) (69.37,0.5515) (69.51,0.5550) (69.64,0.5585) (69.78,0.5622) (69.92,0.5659) (70.05,0.5697) (70.19,0.5735) (70.33,0.5775) (70.47,0.5815) (70.60,0.5856) (70.74,0.5897) (70.88,0.5940) (71.02,0.5983) (71.15,0.6027) (71.29,0.6072) (71.43,0.6118) (71.56,0.6164) (71.70,0.6211) (71.84,0.6259) (71.98,0.6307) (72.11,0.6357) (72.25,0.6407) (72.39,0.6458) (72.52,0.6510) (72.66,0.6562) (72.80,0.6615) (72.94,0.6669) (73.07,0.6724) (73.21,0.6780) (73.35,0.6836) (73.48,0.6893) (73.62,0.6951) (73.76,0.7010) (73.90,0.7069) (74.03,0.7129) (74.17,0.7190) (74.31,0.7252) (74.44,0.7315) (74.58,0.7378) (74.72,0.7442) (74.86,0.7507) (74.99,0.7573) (75.13,0.7639) (75.27,0.7706) (75.41,0.7774) (75.54,0.7843) (75.68,0.7913) (75.82,0.7983) (75.95,0.8055) (76.09,0.8127) (76.23,0.8199) (76.37,0.8273) (76.50,0.8348) (76.64,0.8423) (76.78,0.8499) (76.91,0.8576) (77.05,0.8653) (77.19,0.8732) (77.33,0.8811) (77.46,0.8891) (77.60,0.8972)};
    \addplot[figR, dashed, domain=60:66.76, samples=2] {0.5};
    \addplot[black!55, dotted, line width=0.6pt] coordinates {(60,0.42) (60,0.5)};
    \addplot[black!55, dotted, line width=0.6pt] coordinates {(66.76,0.42) (66.76,0.5)};
    \addplot[only marks, mark=*, mark size=1.5pt, figR] coordinates {(60,0.5) (66.76,0.5)};
    \addplot[only marks, mark=o, mark size=1.8pt, black] coordinates {(63.35,0.4751)};
    \node[font=\scriptsize, figR, anchor=south] at (axis cs:63.4,0.505) {$E'$};
    \node[font=\scriptsize, anchor=south east] at (axis cs:59.8,0.43) {$\theta_1$};
    \node[font=\scriptsize, anchor=south west] at (axis cs:66.95,0.43) {$\theta_2$};
    \draw[black!55, thin] (axis cs:63.7,0.474) -- (axis cs:69.8,0.474);
    \node[font=\scriptsize, anchor=west, black!70] at (axis cs:69.8,0.474) {steady precession};
  \end{axis}
\end{tikzpicture}
